% =============================================================
% TEMPLATE BÁO CÁO (kèm slide) - HCMUS
% Thông tin môn học, đề tài, nhóm: sửa trong config/info.tex.
% Nội dung báo cáo: viết trong content/. Slide: slides.tex.
% Biên dịch: nút xanh trong VS Code, hoặc latexmk -pdf main.tex
% =============================================================
\documentclass[12pt]{article}

\usepackage{hcmus-report-template}
\usepackage{tikz}
\usepackage{caption}
\usepackage{booktabs}
\usepackage{longtable}
\usepackage{tabularx}
\usepackage{array}
\captionsetup{justification=centering}
\setlength{\textfloatsep}{10pt}
\setlength{\floatsep}{10pt}
\usetikzlibrary{shapes, arrows, positioning, matrix, decorations.pathreplacing}

\usepackage[algo2e,ruled,vlined]{algorithm2e}
\DontPrintSemicolon
\SetAlFnt{\small}
\SetKwComment{tcp}{// }{}
\SetAlgorithmName{Thuật toán}{thuật toán}{Danh sách thuật toán}
\SetKwInput{KwIn}{Đầu vào}
\SetKwInput{KwOut}{Đầu ra}
\renewcommand{\lstlistingname}{Mã nguồn}
\usepackage{amsmath}

\setlength{\parindent}{0pt}
\renewcommand{\baselinestretch}{1.5}   % giãn dòng (1 = dòng đơn)
\setlength{\hfuzz}{20pt}
\emergencystretch=3em

% =============================================================
% THÔNG TIN BÁO CÁO - dùng chung cho main.tex (báo cáo) và slides.tex (slide).
% Chỉ cần sửa file này, không phải sửa tên/đề tài ở nhiều nơi.
% =============================================================

% ---------------- Môn học, đề tài ----------------
\newcommand{\coursename}{TÊN MÔN HỌC}              % viết hoa
\newcommand{\reporttitle}{Đồ án số X}               % loại báo cáo: đồ án, bài tập, thực hành...
\newcommand{\reportname}{Tên đề tài}
\newcommand{\reportdate}{\today}

% ---------------- Nhóm thực hiện ----------------
% Mỗi người một dòng, ngăn cách bằng \\ (hiện ở trang bìa)
\newcommand{\studentlist}{Họ và tên 1 --- MSSV --- Lớp\\Họ và tên 2 --- MSSV --- Lớp}
% Bản ngắn cho slide
\newcommand{\studentshort}{Họ và tên 1; Họ và tên 2}
\newcommand{\teachername}{Thầy/Cô: Họ và tên giảng viên}

% ---------------- Trường ----------------
\newcommand{\nationaluni}{Đại học Quốc gia Thành phố Hồ Chí Minh}
\newcommand{\university}{Trường Đại học Khoa học tự nhiên}
\newcommand{\faculty}{Khoa Công nghệ thông tin}

\AtBeginDocument{\renewcommand{\refname}{Tài liệu tham khảo}}

\lhead{\reporttitle: \reportname}
\rhead{
\university\ -- ĐHQG HCM\\
\coursename
}

\begin{document}

% ---- Bìa, mục lục ----
\hypersetup{pageanchor=false}
% Trang bìa. Thông tin lấy từ config/info.tex.
\begin{titlepage}
\newcommand{\HRule}{\rule{\linewidth}{0.5mm}}
\linespread{1}\selectfont   % bìa dùng giãn dòng đơn để vừa một trang
\centering

{\LARGE \MakeUppercase{\nationaluni}}\\[0.5cm]
{\Large \MakeUppercase{\university}}\\[0.5cm]
{\large \MakeUppercase{\faculty}}\\[0.5cm]

\HRule \\[0.4cm]
{
\huge{\bfseries{\reporttitle}}\\[0.5cm]
\large{\bfseries{Đề tài: \reportname}}
}\\[0.4cm]
\HRule \\[0.5cm]

\textbf{\large Môn học: \coursename}\\[0.7cm]

{\large
\begin{tabular}{rl}
\textbf{Sinh viên thực hiện:} & \begin{tabular}[t]{@{}l@{}}\studentlist\end{tabular} \\[0.3cm]
\textbf{Giảng viên hướng dẫn:} & \teachername \\
\end{tabular}}\\[1cm]

{\large \reportdate}\\[1cm]

\IfFileExists{images/hcmus-logo.png}{\includegraphics[scale=.35]{images/hcmus-logo.png}\\[1cm]}{\vspace{1cm}}

\vfill
\end{titlepage}


\pagenumbering{roman}
{\linespread{1}\selectfont\tableofcontents}
\clearpage

\hypersetup{pageanchor=true}
\pagenumbering{arabic}

% ---- Nội dung (đánh số mục) ----
% Muốn thêm chương: tạo file mới trong content/ rồi thêm một dòng \input tương ứng.
% !TEX root = ../main.tex
\section*{Thông tin chung}
\addcontentsline{toc}{section}{Thông tin chung}

\begin{table}[H]
\centering
\caption*{Thông tin thành viên}
\renewcommand{\arraystretch}{1.25}
\begin{tabular}{l c c}
\toprule
\textbf{Thành viên} & \textbf{MSSV} & \textbf{Đóng góp} \\
\midrule
Họ và tên 1 & MSSV & 100\% \\
Họ và tên 2 & MSSV & 100\% \\
\bottomrule
\end{tabular}
\end{table}

\subsection*{Danh mục các phần đã thực hiện}
\begin{itemize}
    \item Phần 1.
    \item Phần 2.
\end{itemize}

\subsection*{Cấu trúc thư mục nộp bài}
\begin{verbatim}
MSSV
|-- baocao.pdf
|-- ...
\end{verbatim}

% !TEX root = ../main.tex
\section{Mở đầu}
Nội dung.

% !TEX root = ../main.tex
\section{Cơ sở lý thuyết}
Nội dung.

% !TEX root = ../main.tex
\section{Nội dung thực hiện}
Nội dung.

% !TEX root = ../main.tex
\section{Kết luận}
Nội dung.


% ---- Tài liệu tham khảo ----
\clearpage
% !TEX root = ../main.tex
% Tài liệu tham khảo lấy từ references.bib. Trích dẫn trong văn bản bằng \cite{khoa}.
% Dòng \nocite dưới đây chỉ để template không báo lỗi khi chưa có \cite nào;
% xoá nó khi báo cáo đã trích dẫn tài liệu thật.
\nocite{example}
\addcontentsline{toc}{section}{Tài liệu tham khảo}
\bibliographystyle{unsrt}
\bibliography{references}


\end{document}
